\documentclass[11pt]{article}

\usepackage[margin=1in]{geometry}
\usepackage[T1]{fontenc}
\usepackage[utf8]{inputenc}
\usepackage{amsmath}
\usepackage{amssymb}

\title{XPC Package Summary\\A Universal Monte Carlo Shapley Toolkit}
\author{Generated technical summary}
\date{\today}

\newcommand{\code}[1]{\texttt{#1}}
\newcommand{\shape}[1]{\ensuremath{\left(#1\right)}}
\newcommand{\flowbox}[1]{\fbox{\begin{minipage}{0.22\linewidth}\centering #1\end{minipage}}}

\begin{document}
\maketitle

\begin{abstract}
XPC is a compact Python package for Monte Carlo Shapley explanations over
explicit feature groups.  It supports ordinary tabular data and temporal tensors
of shape \(\shape{n,d,f}\), where each sample--time location \((n,d)\) can be
explained while the model still receives the complete sequence context.  The
package intentionally exposes one explanation family: Monte Carlo Shapley over
feature-group players.  Kernel SHAP, permutation SHAP, GAM-specific logic,
SMACH-specific logic, generative conditioning, and estimator-side known
contribution handling are outside the estimator.
\end{abstract}

\section{Design Goals}

The package is organized around a small set of explicit abstractions:
\[
\text{\code{DataSpec}},\quad
\text{\code{ModelAdapter}},\quad
\text{\code{Masker}},\quad
\text{\code{Conditioner}},\quad
\text{\code{FeatureGroups}},\quad
\text{\code{ShapleyExplainer}},\quad
\text{\code{Explanation}}.
\]
Each abstraction owns one concern.  Data layout is separate from model calling;
model calling is separate from masking; masking is separate from conditioning;
and signed Shapley estimation is separate from positive post-processing.

\section{End-to-End Flow}

\begin{figure}[h]
\centering
\setlength{\tabcolsep}{4pt}
\begin{tabular}{ccccc}
\flowbox{Input data\\\(\shape{n,f}\) or \(\shape{n,d,f}\)}
& \(\Longrightarrow\) &
\flowbox{\code{DataSpec}\\normalize units}
& \(\Longrightarrow\) &
\flowbox{\code{Masker}\\coalition variants}
\\[6mm]
\flowbox{\code{FeatureGroups}\\players \(G\)}
& \(\nearrow\) & &
\(\searrow\) &
\flowbox{\code{ModelAdapter}\\predictions}
\\[6mm]
& & &
\(\Longrightarrow\) &
\flowbox{\code{ShapleyExplainer}\\Monte Carlo deltas}
\\[6mm]
& & & &
\(\Downarrow\)
\\[2mm]
& & &
\flowbox{\code{Explanation}\\signed values}
& \(\Longrightarrow\)
\flowbox{Heightening\\comparison}
\end{tabular}
\caption{Main runtime flow.  The estimator receives normalized data units,
resolved feature-group players, masked coalition variants, and model
predictions, then returns signed group Shapley values.}
\end{figure}

\section{Data Layout}

\subsection{Tabular Data}

\code{DataSpec} represents independent rows:
\[
X \in \mathbb{R}^{n \times f}.
\]
The explained units are row indices \(i \in \{1,\dots,n\}\).  Model outputs are
normalized to
\[
\hat{y}(X) \in \mathbb{R}^{n \times H},
\]
where \(H=1\) for scalar predictions and \(H>1\) for multi-output models.
The resulting explanation has shape
\[
\Phi \in \mathbb{R}^{n \times H \times G},
\]
where \(G\) is the number of feature-group players.

\subsection{Temporal Tensor Data}

\code{TimeSeriesTensorSpec} represents sequences:
\[
X \in \mathbb{R}^{n \times d \times f}.
\]
The explained units are pairs \((i,t)\).  For a perturbed feature row at time
\(t\), XPC duplicates the complete sequence \(X_i\), replaces only the selected
time-location row, runs the sequence model, and extracts the output at time
\(t\).  Model outputs are normalized to
\[
\hat{y}(X) \in \mathbb{R}^{n \times d \times H},
\]
and explanations have shape
\[
\Phi \in \mathbb{R}^{n \times d \times H \times G}.
\]
This lets temporal models use sequence context without forcing the user to
flatten time into independent rows.

\section{Model Adapters}

\code{ModelAdapter} gives the explainer a uniform prediction boundary:
\[
\code{predict}(X_{\mathrm{batch}}) \rightarrow \mathrm{numpy\ array}.
\]
The package includes adapters for:
\begin{itemize}
  \item NumPy callables and objects with a \code{predict} method;
  \item pandas DataFrame models through \code{PandasModelAdapter};
  \item sklearn-style estimators through \code{SklearnModelAdapter};
  \item PyTorch modules through \code{TorchModelAdapter};
  \item external scripts through \code{ScriptModelAdapter};
  \item R scripts through \code{RScriptModelAdapter};
  \item custom registered classes through \code{register\_model\_adapter}.
\end{itemize}
Adapters do not change the Shapley algorithm.  They only translate data and
prediction outputs.

\section{Feature Groups}

Feature groups are the Shapley players.  A group \(g\) is a disjoint subset of
raw feature indices:
\[
P_g \subseteq \{1,\dots,f\}.
\]
The package supports:
\begin{itemize}
  \item every feature as an individual player;
  \item custom named groups;
  \item remaining features as individual players;
  \item remaining features as one named group;
  \item always-present features.
\end{itemize}
Always-present features are included in every coalition and therefore influence
the base value rather than receiving their own Shapley value.

\section{Masking And Conditioning}

For a unit feature row \(x\) and a present feature set \(S\), a masker creates
masked samples \(\tilde{x}^{(m)}_S\) by preserving \(x_S\) and replacing missing
features \(\bar{S}\).

\subsection{Baseline Masking}

\code{BaselineMasker} replaces missing features with a deterministic baseline.
Accepted baselines are:
\[
\text{zero},\quad \text{mean},\quad \text{median},\quad
\text{scalar},\quad \text{feature vector},\quad
\text{tensor},\quad \text{callable}.
\]
Tensor baselines may be full data-shaped tensors, which is especially useful
for temporal explanations.

\subsection{Random Empirical Masking}

\code{RandomMasker} samples replacement rows from a reference/background matrix
without conditioning:
\[
\tilde{x}_{\bar{S}} \sim \widehat{P}(X_{\bar{S}}).
\]

\subsection{Conditional Masking}

\code{ConditionalMasker} is non-generative.  It delegates to a conditioner that
selects observed reference rows:
\[
\tilde{x}_{\bar{S}} \sim \widehat{P}(X_{\bar{S}} \mid X_S \approx x_S).
\]
The included conditioners are:
\begin{itemize}
  \item \code{EmpiricalConditioner}: nearest empirical neighbors in observed
        feature space;
  \item \code{GridConditioner}: empirical rows in matching grid cells, with
        nearest-neighbor fallback;
  \item \code{CallableConditioner}: user-defined empirical selection.
\end{itemize}

\section{Monte Carlo Shapley Estimator}

For each explained unit \(u\), output horizon \(h\), and feature-group player
\(g\), the Shapley value is
\[
\phi_{u,h,g}
=
\mathbb{E}_{S \subseteq G \setminus \{g\}}
\left[
v_{u,h}(S \cup \{g\}) - v_{u,h}(S)
\right],
\]
where coalitions \(S\) are sampled from the Shapley subset distribution: first
draw a coalition size uniformly from \(0,\dots,G-1\), then draw that many
players uniformly from \(G \setminus \{g\}\).

The coalition value is an empirical model expectation:
\[
v_{u,h}(S)
=
\frac{1}{M}
\sum_{m=1}^{M}
f_h\!\left(\tilde{x}^{(m)}_S\right),
\]
where \(M=\code{n\_mask\_samples}\).  The estimator averages over
\(\code{n\_coalitions}\) sampled coalitions per player.

\begin{figure}[h]
\centering
\begin{tabular}{c}
\fbox{\begin{minipage}{0.55\linewidth}\centering Explained unit \(u\)\end{minipage}}
\\[2mm]\(\Downarrow\)\\[2mm]
\fbox{\begin{minipage}{0.55\linewidth}\centering Player \(g\)\end{minipage}}
\\[2mm]\(\Downarrow\)\\[2mm]
\fbox{\begin{minipage}{0.55\linewidth}\centering Sample \(S \subseteq G\setminus\{g\}\)\end{minipage}}
\\[3mm]
\begin{tabular}{ccc}
\fbox{\begin{minipage}{0.28\linewidth}\centering Mask \(S\)\\predict \(v(S)\)\end{minipage}}
& \(\qquad\) &
\fbox{\begin{minipage}{0.28\linewidth}\centering Mask \(S\cup\{g\}\)\\predict \(v(S\cup g)\)\end{minipage}}
\end{tabular}
\\[3mm]\(\Downarrow\)\\[2mm]
\fbox{\begin{minipage}{0.75\linewidth}\centering
Average deltas:
\(\displaystyle \phi_g \approx \frac{1}{K}\sum_k
\left[v(S_k\cup g)-v(S_k)\right]\)
\end{minipage}}
\end{tabular}
\caption{Estimator loop for one unit and one player.}
\end{figure}

\section{Explanation Object}

\code{Explanation} stores the signed estimator result:
\[
\Phi,\quad b,\quad \hat{y},
\]
where \(\Phi\) has shape \(\mathrm{batch}+ (H,G)\), \(b\) is the base value
with shape \(\mathrm{batch}+(H)\), and \(\hat{y}\) is the model prediction in
the same output shape.  The efficiency residual is
\[
r = \hat{y} - \left(b + \sum_{g=1}^{G} \phi_g\right).
\]
For deterministic baseline masking and additive models this can be exactly
zero; for empirical and conditional Monte Carlo masking it is a diagnostic of
sampling noise and the chosen conditioning distribution.

\section{Precomputed Contributions}

Known or externally supplied contributions are deliberately not used by
\code{ShapleyExplainer}.  Instead they enter through
\[
\code{Explanation.from\_contributions}(\Phi_{\mathrm{ref}}, \dots).
\]
This keeps the estimator mathematically clean while still allowing users to
compare against synthetic ground truth, domain decompositions, or other
explainers.

\section{Heightening}

Heightening is optional post-processing and is enabled by default.  It preserves
signed raw values, shifts them to non-negative values, normalizes to
percentages, and allocates positive target totals:
\[
o_{h,g}
=
\alpha \max\!\left(-\min_u \Phi_{u,h,g}, 0\right),
\qquad \alpha \ge 0,
\]
\[
\Phi^{+}_{u,h,g}
=
\max\!\left(\Phi_{u,h,g} + o_{h,g}, 0\right),
\]
\[
p_{u,h,g}
=
\frac{\Phi^{+}_{u,h,g}}{\sum_j \Phi^{+}_{u,h,j}},
\qquad
\mathrm{part}_{u,h,g}
=
p_{u,h,g} \cdot T_{u,h}.
\]
The target \(T_{u,h}\) defaults to the model prediction and must be
non-negative.  If all positive values are zero, XPC uses a uniform allocation
over groups.  Pass \code{heighten=False} to keep only signed values.

\section{Reference Comparison}

\code{Explanation.compare(reference)} computes aggregate comparison metrics:
mean absolute error, root mean square error, bias, maximum absolute error,
correlation, and per-group mean absolute error.  References may be raw arrays
or another \code{Explanation}; single-output batch vectors are broadcast to the
package's canonical output axis.

\section{Minimal Usage}

\begin{verbatim}
import numpy as np
from xpc import BaselineMasker, ShapleyExplainer

X = np.array([[1.0, 2.0], [3.0, 4.0]])
model = lambda x: x[:, 0] + 2.0 * x[:, 1]

exp = ShapleyExplainer(
    model,
    BaselineMasker("mean"),
    n_coalitions=128,
    random_state=0,
)(X)

print(exp.values.shape)      # (n, H, G)
print(exp.efficiency_residual)
\end{verbatim}

\section{Practical Guidance}

\begin{itemize}
  \item Choose \code{FeatureGroups} according to the explanation question:
        groups are the players, not merely display labels.
  \item Use baseline masking for deterministic counterfactual reference values.
  \item Use random masking for unconditional empirical replacement.
  \item Use conditional masking when missing features should be replaced by
        observed rows compatible with the present coalition.
  \item Increase \code{n\_coalitions} to reduce Shapley subset noise and
        \code{n\_mask\_samples} to reduce masking expectation noise.
  \item Inspect \code{efficiency\_residual}; it is a useful sanity check for
        sampling quality and model-output shape handling.
  \item Keep signed \code{Explanation.values} for analysis.  Heightened parts
        are attached by default for positive target decompositions and can be
        disabled with \code{heighten=False}.
\end{itemize}

\end{document}
